\subsection{Data Preprocessing}
To bridge the gap between raw industrial sensor data and the high-level semantic representations discussed in Section 3.3, a rigorous data preprocessing pipeline is essential. Given that the CLIP-based text encoder operates on refined semantic embeddings in $\mathbb{R}^d$, the raw vibration signals must be transformed into a format that preserves both temporal dependencies and spectral characteristics while being conducive to cross-modal alignment.

The preprocessing begins with raw vibration signal segmentation using a sliding window strategy. Let $X_{raw} \in \mathbb{R}^L$ be a long-duration vibration signal. We apply a sliding window of length $T$ with an overlap ratio $\rho$ to generate a series of samples $\{x_i\}_{i=1}^N$, where each $x_i \in \mathbb{R}^T$. This segmentation ensures that each sample contains sufficient periodic information to capture signature fault frequencies while providing a large volume of training data required for deep manifold learning. The choice of $T$ is typically aligned with the fundamental rotational frequency of the mechanical system to ensure at least 5-10 complete cycles are captured per window.

To extract discriminative features from these non-stationary signals, we employ a Time-Frequency Transformation, specifically the Short-Time Fourier Transform (STFT) or Continuous Wavelet Transform (CWT). By mapping the 1D time-series $x_i$ into a 2D time-frequency representation $S_i \in \mathbb{R}^{F \times T'}$, we can simultaneously observe the evolution of frequency components. This is particularly crucial for detecting transient impacts caused by localized defects in bearings or gears. The resulting spectrograms provide a "visual" analog of the vibration data, which aligns conceptually with the multi-modal nature of the CLIP framework.

Following the transformation, feature normalization and standardization are applied to ensure numerical stability during the diffusion process. Each time-frequency map is scaled such that $S'_i = (S_i - \mu) / \sigma$, where $\mu$ and $\sigma$ are the mean and standard deviation calculated across the "seen" training classes. This prevents dominant frequency components from over-biasing the weight updates in the subsequent Dual-Path Semantic Diffusion module.

Finally, we implement a strict train/test split strategy tailored for the zero-shot diagnosis task. The dataset is partitioned by class labels: $\mathcal{C}_{seen}$ for training and $\mathcal{C}_{unseen}$ for evaluation. Only samples $x \in \mathcal{C}_{seen}$ are used during the initial feature extraction and semantic alignment phase. The model's ability to generalize is then tested on $\mathcal{C}_{unseen}$, where the network must map the preprocessed signals to the latent semantic vectors $s_c$ of previously unencountered fault modes, thereby validating the robustness of the cross-modal knowledge transfer.